\chapter{Objectives}
\label{ch:objectives}

The overarching aim of this research is to study emergent social behaviour,
cooperation, and the value of coordination in Multi-Agent Reinforcement Learning,
using traffic signal control as a controlled, inherently cooperative case study.
The realized objectives of this dissertation are:

\begin{enumerate}[leftmargin=1.6em]
  \item \textbf{Establish a strong cooperative baseline.} Implement and train a
  MAPPO traffic-signal controller on a \grid{2} grid and demonstrate stable
  convergence, then validate it against the standard heuristic controllers
  (fixed-time and max-pressure) under a \emph{fair} evaluation protocol.

  \item \textbf{Confirm implementation fidelity.} Compare the project's MAPPO
  configuration against the canonical published recipe of
  Yu~et~al.~\cite{yu2021mappo}, to establish that the cooperative agent is a faithful,
  competitive implementation and not a weakened strawman.

  \item \textbf{Quantify the value of coordination.} Implement Independent PPO
  (IPPO) as a fully decentralized comparator that differs from MAPPO only in its
  two defining design choices --- a decentralized critic and a purely local reward
  --- and measure the \emph{coordination value}, defined as the performance
  difference $\text{MAPPO} - \text{IPPO}$.

  \item \textbf{Characterize and interpret the result.} Analyze \emph{why} the
  measured coordination value takes the value it does at this network scale, and
  connect the finding to the broader question of how coordination mechanisms scale
  with network complexity.
\end{enumerate}

Two further aims --- characterizing the efficiency--equity trade-off and contrasting
coordination against genuine social dilemmas --- motivate the future work outlined in
Chapter~\ref{ch:futurework}, where network scale-up and a social-dilemma
environment are identified as the natural next steps that compute constraints
placed beyond the scope of this submission.
